\uuid{nZac}
\titre{Pour tout/il existe}
\auteur{bodin ; hanani ; mzari}
\organisation{exo7}
\licence{CC BY-NC-SA 4.0}
\datecreate{2020-01-20}
\niveau{L1}
\module{Algèbre}
\chapitre{Logique, ensemble, raisonnement}
\sousChapitre{Logique}
\difficulte{2}
\theme{Logique, pour tout, il existe}
\exo7qcmid{lille-1.logique.011}

\contenu{
\begin{question}
\qclassification{100.01}
\qtags{motcle=pour tout/il existe}
Soit $f : ]0,+\infty[ \to \Rr$ la fonction définie par $f(x) = \frac{1}{x}$. Quelles sont les assertions vraies ?
\begin{answers}
    \good{$\forall x \in ]0,+\infty[ \quad \exists y \in \Rr \qquad y = f(x)$}

    \bad{$\exists x \in ]0,+\infty[ \quad \forall y \in \Rr \qquad y = f(x)$}

    \good{$\exists x \in ]0,+\infty[ \quad \exists y \in \Rr \qquad y = f(x)$}

    \bad{$\forall x \in ]0,+\infty[ \quad \forall y \in \Rr \qquad y = f(x)$}
\end{answers}
\begin{explanations}
L'ordre des "pour tout" et "il existe" est très important.
\end{explanations}
\end{question}
}
